% from ega2-1-fr.tex
\oldpage[II]{152}
élément transcendant sur $k$. Or, cette vérification est immédiate sur
(\hyperref[II.4.2.6-fr]{4.2.6, exemple 1\textsuperscript{o}}).

\begin{corollary}[7.4.16]
\label{II.7.4.16-fr}
Soient $X$, $Y$ deux courbes algébriques sur $k$, normales, complètes et
irréductibles~; soient $K=R(X)$, $L=R(Y)$ leurs corps. Il existe une
correspondance biunivoque canonique entre l'ensemble des $k$-isomorphismes
$X\xrightarrow{\sim}Y$ et l'ensemble des $k$-isomorphismes
$L\xrightarrow{\sim}K$.
\end{corollary}

C'est une conséquence évidente de (\hyperref[II.7.4.13-fr]{7.4.13}).

\begin{env}[7.4.17]
\label{II.7.4.17-fr}
Le corollaire (\hyperref[II.7.4.16-fr]{7.4.16}) montre qu'une courbe algébrique
sur $k$, normale, complète et irréductible, est \emph{déterminée par son corps
de fonctions rationnelles $K$ à un isomorphisme unique près}~; par définition,
$K$ est une extension de type fini de $k$, de degré de transcendance $1$ (ce
qu'on appelle classiquement un \emph{corps de fonctions algébriques d'une
variable}). De plus~:
\end{env}

\begin{proposition}[7.4.18]
\label{II.7.4.18-fr}
Pour toute extension $K$ de $k$, de type fini et de degré de transcendance
$1$, il existe une courbe algébrique normale, complète et irréductible $X$
telle que $R(X)=K$ (déterminée à isomorphisme unique près). L'ensemble des
anneaux locaux de $X$ s'identifie (\hyperref[I.8.2.1-fr]{I, 8.2.1}) à
l'ensemble formé de $K$ et des anneaux de valuation contenant $k$ et ayant
$K$ comme corps des fractions.
\end{proposition}

En effet, $K$ est une extension de degré fini d'une extension transcendante
pure $k(T)$ de $k$, qui s'identifie, comme on l'a vu, au corps des fonctions
rationnelles de la droite projective $Y=\mathbf{P}_k^1$. Soit $X$ la
\emph{fermeture intégrale} de $Y$ relativement à $K$
(\hyperref[II.6.3.4-fr]{6.3.4})~; $X$ est une courbe algébrique normale de
corps $K$ (\hyperref[II.6.3.7-fr]{6.3.7}), et elle est complète puisque le
morphisme $X\to Y$ est fini (\hyperref[II.7.4.6-fr]{7.4.6}). Les anneaux
locaux $\mathcal{O}_x$ de $X$ sont~: le corps $K$ lorsque $x$ est le point
générique~; si $x$ est distinct du point générique, $\mathcal{O}_x$ est un
anneau de valuation discrète contenant $k$ et ayant $K$ comme corps des
fractions (\hyperref[II.7.4.5-fr]{7.4.5}). Inversement, soit $A$ un tel
anneau~; comme le morphisme $X\to\operatorname{Spec}(k)$ est propre et que
$A$ domine $k$, il domine un anneau local $\mathcal{O}_x$ de $X$
(\hyperref[II.7.3.10-fr]{7.3.10})~; ce dernier étant un anneau de valuation
ayant $K$ comme corps des fractions, est donc nécessairement égal à $A$.

\begin{namedtheorem}[7.4.19]{Remarques}
\label{II.7.4.19-fr}
Il résulte de (\hyperref[II.7.4.16-fr]{7.4.16}) et
(\hyperref[II.7.4.18-fr]{7.4.18}) que \emph{la donnée d'une courbe algébrique
sur $k$, normale, complète et irréductible, est essentiellement équivalente à
la donnée d'une extension $K$ de $k$, de type fini et de degré de
transcendance $1$}. On notera que si $k'$ est une extension du corps de base
$k$, $X\otimes_k k'$ sera encore une courbe algébrique complète sur $k'$
(\hyperref[II.5.4.2-fr]{5.4.2, (iii)}), mais en général elle ne sera ni réduite
ni irréductible. Elle le sera cependant si $K$ est une extension séparable de
$k$ et si $k$ est algébriquement fermé dans $K$ (ce qui s'exprime, dans une
terminologie classique que nous ne suivrons pas, en disant que $K$ est une
«~extension régulière~» de $k$). Mais même dans ce cas, il peut se faire que
$X\otimes_k k'$ ne soit pas normale. Le lecteur trouvera des détails sur ces
questions dans le chapitre IV.
\end{namedtheorem}

\section{Schémas éclatés~; cônes projetants~; fermeture projective}
\label{section:II.8-fr}

\subsection{Préschémas éclatés}
\label{subsection:II.8.1-fr}

\begin{env}[8.1.1]
\label{II.8.1.1-fr}
Soient $Y$ un préschéma, et pour tout entier $n\geq0$, soit $\mathcal{I}_n$ un
Idéal quasi-cohérent de $\mathcal{O}_Y$~; on suppose vérifiées les conditions
suivantes~:
\begin{equation}
\label{II.8.1.1.1-fr}
  \mathcal{I}_0=\mathcal{O}_Y,\qquad
  \mathcal{I}_n\subset\mathcal{I}_m\quad\text{pour }m\leq n
\tag{8.1.1.1}
\end{equation}
\begin{equation}
\label{II.8.1.1.2-fr}
  \mathcal{I}_m\mathcal{I}_n\subset\mathcal{I}_{m+n}
  \quad\text{quels que soient }m,n.
\tag{8.1.1.2}
\end{equation}

